\section{Results}
\label{sec:results}

\subsection{Controlled dilution and optimizer diagnosis}
\label{sec:optimizer}

The scale contract in \cref{fig:dilution-scale} removed the large numerical-
span confound while preserving all exact optima. The nested results are
summarized in \cref{fig:optimizer-attribution,tab:optimizer-attribution}. The 56-task depth pilot then
completed 504/504 planned optimizations. Its independently initialized $p=3$
fit appeared worse than $p=2$, but the nesting test changed that
interpretation: all 168 zero-layer identities held exactly, and 29/168
reported $p=3$ mean-energy objectives were worse than their embedded $p=2$
points. Continuation lowered the objective in 29/29 certified failures and
increased $\gfeas$ in 27/29. The remaining 139 runs simply passed this
necessary test; they were not certified global optima.

\begin{figure}[t]
  \centering
  \includegraphics[width=\textwidth]{fig03_optimizer_attribution.pdf}
  \caption{Nested optimizer diagnosis. \textbf{a}, original $p=3$ objective
  relative to the exactly embedded $p=2$ point; 29 positive gaps certify
  failure. \textbf{b}, matched feasibility gain before and after continuation}
  \label{fig:optimizer-attribution}
\end{figure}


\begin{table}[t]
\centering
\caption{Optimizer attribution on the 56-task discovery set. Counts are seed-level unless stated otherwise.}
\label{tab:optimizer-attribution}
\small
\begin{tabularx}{\textwidth}{@{}Xr@{}}
\toprule
Diagnostic & Result \\
\midrule
Nested $p=2\to p=3$ zero-layer identity & 168/168 \\
Original $p=3$ objective worse than embedded $p=2$ & 29/168 \\
Continuation recovered lower objective in that failure set & 29/29 \\
Continuation increased $G_{\mathrm{feas}}$ in that failure set & 27/29 \\
Random $p=3$ had lower energy but lower $G_{\mathrm{feas}}$ & 82/168 \\
\bottomrule
\end{tabularx}
\end{table}


Mean-energy improvement was not equivalent to feasible-mass improvement. In
82/168 matched comparisons, the state with lower mean energy also had lower
$\gfeas$; component accounting showed redistribution among infeasible energy
levels. Certified optimizer inadequacy also occurred under every tested local
optimizer. O0/O3 nested-failure counts were 28/168 and 38/168 for COBYLA,
54/168 and 61/168 for Nelder--Mead, and 46/168 and 72/168 for SLSQP.
Accordingly, certified optimizer inadequacy occurs under several tested
optimizers and evaluation budgets; this is not a claim about arbitrary
classical budgets or extensive multi-start search.

\subsection{Capacity control and objective alignment}
\label{sec:objective-alignment}
\label{sec:heldout}

On discovery, O2 increased $\pfeas$ over O0 on 53/56 tasks. The median
O2--O0 capacity gap was 0.2193 decades in $\gfeas$, demonstrating feasible-
concentration capacity that mean-energy training did not induce under the
matched design. O3 increased $\pfeas$ over O0 on 50/56 tasks, left a median
O2--O3 residual of 0.00690 decades, and closed a median 97.9\% of the
taskwise O2--O0 gap. These discovery results selected O3 and
$\alpha=0.10$; O2 remains an exact-statevector capacity control, not a
deployable objective.

Both preregistered held-out criteria passed
(\cref{fig:heldout,tab:heldout-tests}). Across 15 held-out graph units and 84 tasks, H1 estimated
O3--O0 $\gfeas$ at $+0.3547$ decades, with one-sided 95\% lower bound
$+0.2374$ and Holm-adjusted $p=2.44\times10^{-4}$. H2 estimated O3--O2 at
$-0.0086$ decades; its lower bound $-0.0360$ exceeded the frozen
$-0.10$-decade non-inferiority margin, with the same adjusted $p$-value. Thus
CVaR approached the measured exact-feasibility capacity at the graph analysis
unit. O3 also increased $\popt$ on 80/84 tasks, while changes in $\pcond$
remained heterogeneous.

\begin{figure}[t]
  \centering
  \includegraphics[width=\textwidth]{fig06_heldout_confirmation.pdf}
  \caption{Preregistered held-out objective alignment. Points are complete
  graph means; solid lines are mean effects, dotted lines are one-sided 95\%
  lower bounds, and dashed lines are the H1 null or H2 non-inferiority margin}
  \label{fig:heldout}
\end{figure}


\begin{table}[t]
\centering
\caption{Preregistered held-out graph-level tests ($15$ base graphs; $84$ tasks). Lower bounds are one-sided grouped-bootstrap 95\% bounds; $p$-values are Holm adjusted.}
\label{tab:heldout-tests}
\small
\begin{tabularx}{\textwidth}{@{}lXrrrr@{}}
\toprule
Test & Contrast in $G_{\mathrm{feas}}$ & Effect & Lower bound & Null/margin & Holm $p$ \\
\midrule
H1 & O3 CVaR $-$ O0 mean energy & $+0.3547$ & $+0.2374$ & $0$ & $0.000244$ \\
H2 & O3 CVaR $-$ O2 capacity control & $-0.0086$ & $-0.0360$ & $-0.10$ & $0.000244$ \\
\bottomrule
\end{tabularx}
\end{table}


The distinction survives exclusion of certified failures. Requiring both O0
and O3 to pass the nested test leaves 113, 72, and 70 pairs under COBYLA,
Nelder--Mead, and SLSQP; their graph-cluster O3--O0 $\gfeas$ intervals are
$[0.5770,0.9971]$, $[0.5796,1.2473]$, and $[0.6096,1.3458]$. This supports
separation of optimizer inadequacy from objective misalignment under the
tested diagnostic, not universal independence.

\subsection{Depth and robustness boundaries}
\label{sec:depth-robustness}

In the preselected 24-task depth--budget subset, pooled $p=2\to3$ nested-
failure rates were 20.8\%, 21.5\%, and 26.4\% at 120, 240, and 480 calls;
$p=3\to4$ rates were 31.2\%, 31.2\%, and 32.6\%. Increasing the tested
evaluation budget did not monotonically reduce nested failures in this
controlled subset. Yet at 120 calls, graph-mean $p=4$ minus $p=3$ $\gfeas$
was $+0.7206$ for O0 and $+0.2021$ for O3, positive on 10/10 graphs in each
arm; the analogous effects remained positive at 240 and 480 calls. Depth can
therefore add usable capacity while making the tested optimization problem
harder. We do not claim that these instances remain difficult under
arbitrarily large classical optimization budgets or extensive multi-start
search.

The post-hoc discovery scan also bounded tail-fraction sensitivity. Equal-
weight graph means were similar around the frozen choice for
$\alpha=0.05$, 0.10, and 0.25, but optimized responses were non-monotone on
54/56 tasks. At $\alpha=1$, exact weighted CVaR equaled mean energy to within
$2.42\times10^{-13}$; separately executed optimizer trajectories need not
coincide. No post-hoc tail fraction revises the frozen H1/H2 result.

\subsection{Finite-shot and scaling limits}
\label{sec:scaling}

Fixed-endpoint sampling recovered the exact sign of non-tied O3--O0
$\pfeas$ comparisons in 94.8\%, 99.3\%, and 100.0\% of cases at
$10^3$, $10^4$, and $10^5$ shots. End-to-end finite-shot training was less
decisive (\cref{fig:robustness-boundaries}). The matched exact-trained effect
on ten held-out graphs was 0.3599
(95\% graph-cluster CI $[0.1641,0.5697]$); effects were 0.1706
($[-0.0106,0.3862]$) at $10^4$ shots per call and 0.1246
($[-0.0050,0.3351]$) at $10^3$. Both point estimates retained the O3--O0
direction, but both intervals crossed zero.

\begin{figure}[t]
  \centering
  \includegraphics[width=\textwidth]{fig12_reviewer_depth_budget.pdf}
  \includegraphics[width=\textwidth]{fig15_reviewer_finite_shot_training.pdf}
  \caption{Robustness boundaries. Upper: nested diagnostics and terminal
  quality across tested depths and evaluation budgets. Lower: exact-terminal
  evaluation after exact or finite-shot training; both shot-trained graph
  intervals cross zero}
  \label{fig:robustness-boundaries}
\end{figure}

The completed size-conditioned study also prevents extrapolation. At $m=20$,
mean response exponents were 0.6150 for O0, 0.7885 for O2, and 1.2821 for O3,
reversing the earlier O3--O0 ordering. At $m=22$, runs exceeded the
preregistered computational budget and were administratively censored. This
is a limitation of the study's computational budget, not a principled quantum-
algorithm barrier, a claim that 22 qubits cannot be simulated, or a
statevector-memory limit. We therefore make no extrapolative scaling claim
beyond the completed range.

Finally, an exact label-setting solver reproduced 140/140 frozen optima with
small label sets and microsecond-scale descriptive timings. The benchmark is
classically tractable by design. Complete optimizer, depth, alpha, finite-shot,
classical, and scaling matrices are in Online Resource~1, Secs.~S4--S8.
